\begin{helper}[Displayed-edge preservation through the last starred stage]
Let $G$ be a finite loopless multigraph with exact edge identities, palette
$[k]$, and distinguished edge $e_0$.  Fix an ETT $\mathcal E$ on $T$
generated by $\varphi$, a \noderef{HierarchyBase} $B$, a
\noderef{HierarchyCounterexample} $C$ on $B$, and a
\noderef{HierarchyMinimumCounterexampleContext} for these data.
Write $H=C.\mathit{hierarchy}$, $q=H.\mathit{dividerCount}$,
$n=\mathcal E.\mathit{rounds}$, $T_n$ for its terminal ladder tree,
$D_n$ for its computed terminal stored set, $R=B.\mathit{closure}$,
$A=\noderef{HierarchyStarredStage}(H,q)$, and
\[
 Y=T|_{\min\{|A|+(|T|-|A|-1),|T|\}}.
\]
Lengths count displayed edges and subtraction is truncated.  Assume
$Y$ elementary under $\varphi$.  Let distinct colors $\alpha,\beta$ be
missing on $Y$, let $Q$ be the full component through a node $z$, with a
\noderef{TwoColorPathComponent} certificate, and let
$\sigma=\noderef{TwoColorComponentSwap}(G,\varphi,\alpha,\beta,z)$.
Put $a=T.\mathit{vertex}(\noderef{HierarchyFirstMissingPosition}(\varphi,T,\alpha))$.
Assume $Q$ disjoint from the full $\alpha,\beta$ component through $a$, and
assume
\[
 (q>0\ \land\ \alpha\in\overline\varphi(H.\mathit{stage}(q)))
 \quad\text{or}\quad
 \bigl(q=0\ \land\
   [\,\alpha\in\overline\varphi(T_n)\ \text{or}\
      (\alpha\in\overline\varphi(A),\ \alpha,\beta\notin D_n)\,]\bigr).
\]
Then every exact displayed edge of $A$, including the first edge $e_0$,
has the same value under $\sigma$ and $\varphi$.
\end{helper}
\begin{proof}
Apply \noderef{P2813_T1_BaseControl} to obtain the base stability relation,
properness and one-uncoloredness of $\sigma$, missing-set equality on
$A$, and outside recoloring on every starred stage carrying an affected
missing closed color.  Apply \noderef{P2813_T1_BaseDisplayedEdges} to fix
all displayed edges of the core.  The exact starred-stage inclusions from
\noderef{TashkinovHierarchy} put the core inside $A$; its edges retain
the same identities in $A$.  If $q=0$, $A$ is the core, so this proves
the claim.  Any edge of $A$ beyond the core therefore has $q>0$.

The predecessor-elementarity premise for \noderef{rutcor} follows from
\noderef{TechnicalTheoremBase} when $n=1$, or from the actual maximum
predecessor supplied by \noderef{CertifiedETTPredecessor} and the context's
lower assertion (i) when $n>1$.  Its prefix-elementarity premise is
the assumed elementarity of $Y$.  Exact-prefix vertex inclusion and
\noderef{ColoredVertexSetElementary} make every starred stage elementary.

Fix an affected color $c$ missing on $A$.
If it has a closed starred stage, take its largest such index $r$.
It is also missing there: otherwise its least first-missing starred
index $s$ is larger than $r$, and
\noderef{GoodHierarchyNewMissingClosed} makes it closed at $s$, a
contradiction.  Missingness at $r$ follows by exact-prefix inclusion if
the least index is at most $r$.  The maximal-closed clause of
\noderef{rutcor} says no displayed edge at or beyond $|A_r|$ through $A$
has color $c$.  Any displayed $c$-edge before that position is an exact
edge of $A_r$, and hence is fixed by
\noderef{OutsideKempeRecoloringStable}, applied to the outside recoloring
returned by \noderef{P2813_T1_BaseControl}.
If $c$ has no closed starred stage, the no-closed clause of
\noderef{rutcor} excludes $c$ from all displayed edges of $A$ at or
beyond $|T_n|$, so in particular from those beyond the core.
Thus every displayed affected edge whose color is missing on $A$
is fixed.

Now fix an affected color $c$ not missing on $A$.  Obtain the actual
terminal-route alternative from \noderef{P2813_T1_TerminalRoute}.
An ordinary augmentation edge at an index $t$ in $A$ cannot have color
$c$: \noderef{TashkinovAugmentationSequencePosition} on the route to
$T$ supplies a missing witness on the exact prefix $T|_t$ for its
color.  Since $t<|A|$, that vertex also belongs to $A$, and
\noderef{TreeMissingColors} would make $c$ missing on $A$.

In PE every edge beyond $T_n$ is such an augmentation edge.
In SE/RE the same argument applies at indices at or beyond the seed
length.  The sole remaining position is $t=|T_n|$, the actual connecting
seed edge $f$: equality of the seed and $A$ entries with those of $T$,
and \noderef{TreeSequenceAdjoinEdge}, identify this exact edge identity
with $f$.  The terminal-route helper states both that $f$ is incident
with $T_n$ and that its old color belongs to
$\overline\varphi(T_n)\cup D_n$.
The $T_n$ stability relation from the base-control helper therefore fixes
$f$ by the incident-edge clause of \noderef{StableColoring}.
This explicitly handles the stored-seed exception.

Finally an edge whose old value is neither $\alpha$ nor $\beta$ is fixed
directly by \noderef{TwoColorComponentSwap}, regardless of whether its
edge node lies in the component.  Together with the core, missing-color,
and connecting-seed cases above, this covers every exact displayed edge.
\end{proof}

